\section{Revisión de contenidos}

A continuación, se presentan una serie de preguntas diseñadas para evaluar la compresión de los conceptos fundamentales presentados en el texto.

\subsection{Preguntas sobre Procesos}

\begin{enumerate}
  \item ¿Qué es un proceso y en qué se diferencia de un programa?
  \item ¿Qué información almacena el Bloque de Control de Proceso (PCB) y por qué es indispensable para el sistema operativo?
  \item ¿Qué es la imagen de un proceso y cuáles son sus componentes principales?
  \item ¿Qué es un cambio de contexto (context switch) y por qué es necesario?
  \item ¿Cuáles son los estados básicos por los que puede pasar un proceso y qué representa cada uno?
  \item Realizar el diagrama de 7 estados de un proceso (5 estados $+$ estado suspendido).
  \item ¿Qué eventos pueden provocar la creación de un proceso?
  \item ¿Cuáles son las causas más comunes de terminación de un proceso?
  \item ¿Qué tablas mantiene el sistema operativo para administrar los recursos del sistema?
  \item ¿Por qué los procesadores implementan modo usuario y modo kernel?
  \item ¿Qué diferencia existe entre un cambio de modo (mode switch) y un cambio de proceso (process switch)?
  \item ¿Cuál es la diferencia entre una interrupción y una excepción o trap?
  \item ¿Qué pasos realiza el sistema operativo para crear un nuevo proceso?
  \item ¿Por qué el PCB es considerado la estructura de datos más importante para la gestión de procesos?
  \item ¿Qué ventajas aporta el modelo de procesos en sistemas multiprogramados?
  \item ¿Cuáles son los principales modelos de ejecución del núcleo y en qué se diferencian?
  \item ¿Por qué en la ejecución del núcleo dentro del proceso de usuario normalmente solo ocurre un mode switch y no un process switch?
\end{enumerate}

\subsection*{Respuestas}

\paragraph{Respuesta 1} 
Un programa es un conjunto de instrucciones almacenadas en disco. Un proceso es un programa en memoria principal que puede ser ejecutado, junto con sus datos, recursos y contexto de ejecución administrados por el sistema operativo.

\paragraph{Respuesta 2}
Almacena el identificador del proceso, su estado, el contexto del procesador (registros y PC), información de planificación, memoria, recursos e información de contabilidad. Permite suspender y reanudar la ejecución de un proceso sin perder su estado.

\paragraph{Respuesta 3}
Es la representación completa de un proceso en memoria. Está formada por:
\begin{itemize}
  \item Código del programa.
  \item Datos.
  \item Pila.
  \item Bloque de Control de Proceso (PCB).
\end{itemize}

\paragraph{Respuesta 4}
Es la operación mediante la cual el sistema operativo guarda el estado del proceso que estaba ejecutándose y restaura el estado de otro proceso para continuar su ejecución.

\paragraph{Respuesta 5}
\begin{itemize}
  \item Nuevo (New): el proceso acaba de ser creado.
  \item Listo (Ready): puede ejecutarse, pero espera el procesador.
  \item Ejecutando (Running): está utilizando el procesador.
  \item Bloqueado (Blocked/Waiting): espera un evento, generalmente una operación de E/S.
  \item Terminado (Exit): finalizó su ejecución.
\end{itemize}

\paragraph{Respuesta 6}
Ver figura \ref{fig:five_state_model_with_suspend}.

\paragraph{Respuesta 7}
\begin{itemize}
  \item Inicio de un trabajo (batch).
  \item Inicio de sesión de un usuario.
  \item Creación por parte del sistema operativo para brindar un servicio.
  \item Creación por otro proceso (spawn).
\end{itemize}

\paragraph{Respuesta 8}
\begin{itemize}
  \item Finalización normal.
  \item Error de ejecución (división por cero, instrucción inválida, etc.).
  \item Violación de memoria o protección.
  \item Intervención del sistema operativo o del usuario.
  \item Terminación del proceso padre.
\end{itemize}

\paragraph{Respuesta 9}
Generalmente mantiene tablas de:
\begin{itemize}
  \item Procesos.
  \item Memoria.
  \item Dispositivos de entrada/salida.
  \item Archivos.
\end{itemize}
Estas estructuras contienen toda la información necesaria para administrar el sistema.

\paragraph{Respuesta 10}
Para proteger el sistema operativo y los recursos críticos. Las instrucciones privilegiadas y el acceso a recursos sensibles solo pueden ejecutarse en modo kernel.

\paragraph{Respuesta 11}
Un cambio de modo modifica únicamente el nivel de privilegio del procesador (usuario \(\leftrightarrow\) kernel), mientras continúa ejecutándose el mismo proceso.

Un cambio de proceso implica guardar el contexto de un proceso y restaurar el de otro, cambiando completamente la ejecución.

\paragraph{Respuesta 12}
Una interrupción es un evento generalmente externo al proceso (por ejemplo, un dispositivo de E/S o el temporizador).

Un trap o excepción es un evento interno provocado por la ejecución de una instrucción, como una llamada al sistema o una división por cero.

\paragraph{Respuesta 13}
\begin{enumerate}
  \item Asigna un identificador único (PID).
  \item Reserva memoria para la imagen del proceso.
  \item Inicializa el PCB.
  \item Inserta el proceso en las estructuras de planificación.
  \item Actualiza las tablas de administración correspondientes.
\end{enumerate}

\paragraph{Respuesta 14}
Porque contiene toda la información necesaria para administrar un proceso. Gracias al PCB el sistema operativo puede planificar procesos, asignar recursos, atender interrupciones y realizar cambios de contexto.

\paragraph{Respuesta 15}
Permite ejecutar múltiples programas de forma concurrente, compartir eficientemente el procesador y otros recursos, aislar procesos entre sí y simplificar la planificación y administración del sistema operativo.

\paragraph{Pregunta 16}
Existen tres modelos principales:
\begin{itemize}
  \item Núcleo independiente de los procesos (Nonprocess Kernel): el kernel se ejecuta como una entidad separada de los procesos de usuario, utilizando su propia memoria y pila.
  \item Ejecución dentro del proceso de usuario: el kernel se ejecuta en el contexto del proceso que solicitó un servicio o fue interrumpido. Normalmente solo ocurre un cambio de modo (mode switch), y solo se realiza un cambio de proceso si el planificador lo decide.
  \item Sistema operativo basado en procesos (Process-Based OS): las funciones del sistema operativo se implementan como procesos independientes, favoreciendo una arquitectura modular y facilitando la ejecución concurrente, especialmente en sistemas multiprocesador.
\end{itemize}

\paragraph{Pregunta 17}
Porque el procesador continúa ejecutando el mismo proceso que originó la interrupción o la llamada al sistema. Únicamente cambia el nivel de privilegio de usuario a kernel (mode switch). Solo si el planificador decide ceder el procesador a otro proceso será necesario realizar un process switch, guardando el contexto del proceso actual y restaurando el de otro.

\subsection{Preguntas sobre hilos}

\begin{enumerate}
  \item ¿Por qué se dice que un proceso es la unidad de asignación de recursos y un hilo la unidad de ejecución?
  \item ¿Por qué cada hilo necesita su propia pila de ejecución, aunque comparta la memoria del proceso?
  \item ¿Qué recursos son compartidos por los hilos de un mismo proceso y cuáles son privados?
  \item ¿Por qué la comunicación entre hilos suele ser más eficiente que la comunicación entre procesos?
  \item ¿Qué problema puede aparecer cuando varios hilos comparten memoria?
  \item ¿Por qué los cambios de contexto entre hilos de usuario (ULT) son más rápidos que entre hilos de kernel (KLT)?
  \item ¿Por qué una llamada al sistema bloqueante puede detener a todos los hilos en una implementación ULT?
  \item ¿Qué ventaja tienen los KLT frente a los ULT en un sistema con varios procesadores?
  \item ¿Cuál es el objetivo de un modelo híbrido que combina ULT y KLT?
  \item ¿En qué situaciones resulta conveniente utilizar múltiples hilos en lugar de múltiples procesos?
\end{enumerate}

\subsection*{Respuestas}

\paragraph{Respuesta 1} Porque el proceso posee el espacio de direcciones y los recursos (memoria, archivos, dispositivos, etc.), mientras que el hilo representa el flujo de ejecución que utiliza dichos recursos.

\paragraph{Respuesta 2} Porque cada hilo mantiene su propio contexto de ejecución, incluyendo llamadas a funciones, variables locales y direcciones de retorno. Si compartieran la misma pila, interferirían entre sí.

\paragraph{Respuesta 3} Comparten el espacio de direcciones, el código, los datos globales, los archivos abiertos y otros recursos del proceso. Cada hilo posee su propio contador de programa, registros, pila y contexto de ejecución.

\paragraph{Respuesta 4} Porque los hilos comparten la misma memoria y pueden intercambiar información directamente, mientras que los procesos requieren mecanismos de comunicación proporcionados por el sistema operativo.

\paragraph{Respuesta 5} Pueden producirse condiciones de carrera e inconsistencias si varios hilos modifican simultáneamente los mismos datos. Por ello es necesario utilizar mecanismos de sincronización.

\paragraph{Respuesta 6} Porque los ULT son administrados completamente en espacio de usuario y no requieren cambiar al modo kernel. En cambio, los KLT necesitan la intervención del sistema operativo.

\paragraph{Respuesta 7} Porque el kernel desconoce la existencia de los hilos y únicamente observa al proceso. Si el proceso se bloquea por una llamada al sistema, ninguno de sus hilos puede continuar ejecutándose.

\paragraph{Respuesta 8} El kernel puede planificar distintos hilos de un mismo proceso sobre diferentes procesadores, permitiendo su ejecución en paralelo.

\paragraph{Respuesta 9} Combinar la baja sobrecarga y flexibilidad de los ULT con la capacidad de paralelismo y manejo de bloqueos que ofrecen los KLT.

\paragraph{Respuesta 10} Cuando varias tareas pertenecen a una misma aplicación y necesitan compartir datos o recursos de forma eficiente, ya que los hilos son más rápidos de crear, sincronizar y cambiar de contexto que los procesos.
